% Figure 1 Caption: Model Preference Heterogeneity (Revised)

\begin{figure}[t]
\centering
\includegraphics[width=0.95\textwidth]{experiments_v1/01_figure/results/figure1_lmsys_holdout_pca.png}
\caption{\textbf{Model Preference Heterogeneity in LMSYS Held-Out Prompts ($N=750$).}
\textbf{(A)} PC1 projection (domain-adapted PCA, trained on 80K routing prompts, unsupervised) vs.\ reward gap $(R_{\text{GPT-4-Turbo}} - R_{\text{Mixtral}})$. Reward gaps are discrete pairwise outcomes. The dashed line marks the silhouette-optimal threshold ($\text{PC1}=0.222$), selected without reference to reward labels. The running mean shows a smooth transition in model preference across the PC1 axis.
\textbf{(B)} Outcome proportions by cluster. The Low PC1 cluster (80.8\%, $n=606$, blue) and High PC1 cluster (19.2\%, $n=144$, red) exhibit significantly different outcome distributions ($\chi^2$ test: $p < 0.0001$; Mann-Whitney $p < 0.0001$). The High PC1 cluster shows a substantially higher Mixtral win rate and lower GPT-4-Turbo win rate. The primary, unbiased effect size from generic C4 PCA (100K web text, no routing connection) is Cohen's $d = 0.33$ (small effect, $p < 0.0001$); domain-adapted PCA amplifies this to $d = 1.53$ via task-relevant feature extraction.
Note: reward gaps are discrete pairwise outcomes (win/tie/loss) from LMSYS Chatbot Arena human preference evaluations; the running mean in Panel~A smooths outcome proportions for visual clarity. Cluster proportions are specific to this holdout sample ($\sim$5.9\% High PC1 in the 1M dataset).}
\label{fig:lmsys_holdout_structure}
\end{figure}
